\section{总结与延伸}

Lecture 9 给出的不是一个公式，而是一套训练前预测方法。先从历史和数据 scaling 建立 power-law 直觉，再用模型工程 scaling 比较架构、优化器和超参数，最后用联合 data-model scaling 决定 compute、parameter、token 的分配。它把大模型训练从“大赌注”变成一组可审计的小实验和外推。

\begin{importantbox}{最终 takeaways}
\begin{enumerate}
    \item Scaling law 常见形式是 \(L(x)=A x^{-\alpha}+C\)，在 log-log 图上近似直线。
    \item 数据 scaling 的 slope 反映继续加数据的边际收益，offset 常受数据 mixture 和分布影响。
    \item 小模型 scaling 实验可以预测架构、optimizer、depth/width、batch size 和 learning rate 的大规模趋势，但必须控制变量。
    \item Critical batch size 随目标 loss 变化，batch 不是越大越好。
    \item muP 的价值是让小模型超参数更可靠迁移到大模型。
    \item Joint scaling law 决定更多数据还是更大模型；Chinchilla 改变了 compute-optimal 训练的直觉。
    \item Train-optimal 不等于 deployment-optimal，推理成本会推动 overtraining。
\end{enumerate}
\end{importantbox}

\subsection{拓展阅读}

阅读 scaling-law 论文时，不要只摘录最终 exponent。至少记录模型与数据网格、参数/FLOPs 口径、optimizer 和 schedule、每条 curve 是否训练充分、fit range、residual、外推误差与目标函数（train loss、downstream metric 或 total deployment cost）。这些信息决定两条看似相同的 power law 能否比较。

\begin{longtable}{L{0.24\textwidth}L{0.30\textwidth}L{0.36\textwidth}}
\toprule
材料 & 为什么值得读 & 带着什么问题读 \\
\midrule
Hestness et al. 2017 & 系统展示多任务 data/compute scaling 的早期证据 & 不同 domain 的 slope/irreducible error 如何变化，emergence 是否只是阈值投影？ \\
Kaplan et al. 2020 & 建立 LLM 参数、数据、compute 与 loss 的经典联合框架 & 参数计数、warmup、run range 与 undertraining 怎样影响 $N_{opt},D_{opt}$？ \\
Hoffmann et al. 2022（Chinchilla） & 给出 IsoFLOPS、joint fit 与更 data-heavy 的 compute-optimal 配比 & 三种 fitting methods 是否一致，method 3 的数据/实现风险是什么？ \\
Rosenfeld et al. 2020 & 讨论跨模型规模预测 generalization error 的构造方法 & 哪些结构可由小模型 extrapolate，哪些常数依赖 architecture/dataset？ \\
Yang et al. 的 $\mu$P 工作 & 让 learning rate、initialization 等超参数更稳定地跨 width transfer & 哪些超参数保持不变，哪些必须按 parameterization 缩放？ \\
部署经济学与 overtraining 分析 & 把训练 compute 扩展成训练+重复推理的 total cost & 在预期请求量、latency 和显存约束下，最优模型为何可能小于 Chinchilla 模型？ \\
\bottomrule
\caption{Lecture 9 的注释式拓展阅读路线。}
\end{longtable}

\begin{importantbox}{建议的复现实验}
训练至少四个模型尺寸和四个 token budgets，保存完整 loss curves。分别复现 data-only power law、固定 FLOPs 的 IsoFLOPS minima 和 joint $L(N,D)$ fit；统一 total/non-embedding parameter 口径，并用一个未参与拟合的更大 run 检验外推。最后把 inference 请求量加入目标，比较 train-optimal 与 deployment-optimal 配置。
\end{importantbox}

\end{document}
